\scanpage{262}
\mspage{74}
% archive sheet number: 262
comme une sorte de (\unsure{pâle}) ``squelette'' du diagramme substantiel dont il provient. P. ex. si

\diagram{\begin{tikzcd}[ampersand replacement=\&,column sep=1.6em]
I=\mathbb{\Delta}_0:\quad 0\arrow[r]\&1
\end{tikzcd}}

on note un diagramme substantiel \inserted{\unsure{de|du} type $(\mathbb{\Delta}_{\unsure{1|0}})^{\unsure{\mathrm o|\circ}}$} par le même graphisme, que le \deleted{\illegible{1}} diagramme ordinaire :

\diagram{\begin{tikzcd}[ampersand replacement=\&] F_0\&F_1\arrow[l]\end{tikzcd}\,.}

Si

\diagram{$I={}$\begin{tikzcd}[ampersand replacement=\&,row sep=small,column sep=large]
a\arrow[r]\&c\\
0\arrow[u]\arrow[r]\&b\arrow[u]
\end{tikzcd}}

les diagrammes substantiels de type $I$ se notent

\diagram{\begin{tikzcd}[ampersand replacement=\&,row sep=small,column sep=large]
F_a\arrow[d]\&F_c\arrow[l]\arrow[d]\\
F_0\&F_b\arrow[l]
\end{tikzcd}\,.}

Il faudra bien sûr préciser chaque fois si ce graphisme désigne un diagramme squelettique, ou substantiel. Souvent on part d'un squelette, et on essaye de ``l'étoffer'' avec un peu de chair, pour en faire un diagramme substantiel --- si possible, de façon canonique.

C'est là un point de philosophie dérivatique \unsure{essentielle|essentiel} pour la pratique. On \unsure{procédera|procèdera} souvent de proche en proche, en ``\unsure{substantifiant}'' des sous-\deleted{\unsure{g}}diagrammes correspondant \unsure{à|:} une suite strictement croissante $I_0\subset I_1\ldots\subset I_n=I$ de $I$, au mieux de la situation particulière.
